\section{Proof of the slope-dependent output-error bound}
\label{app:slope-proof}

The continuous-center integral exposes how the common tanh slope changes the
readout kernel. An explicit lattice allowance transfers this calculation to
finitely many equally spaced centers. Interlacing supplies upper bounds on the
relative learning rates; substituting these bounds into the exact GD residual
formula gives the output-error lower bound. We separate this analytic argument
from its numerical evaluation below.

\paragraph{Setup and statement.}
Fix distinct samples $z_1,\ldots,z_m$ and $W$ consecutive centers $x_j$ on
$c_0+h\mathbb Z$, with $h>0$, covering the sample interval. Fix $\gamma>0$ and set
\begin{equation}
 J_\gamma=\frac1{\sqrt m}
 [\mathbf1,(\tanh(\gamma(z_a-x_j)))_{a,j}],\qquad
 K_\gamma=J_\gamma J_\gamma^T,\qquad y_a=\frac{f(z_a)}{\sqrt m}.
 \label{eq:app-slope-setup}
\end{equation}
Thus $\|J_\gamma w-y\|^2$ is the mean squared sampled output error. Assume
$y\ne0$, write the descending eigenpairs of $K_\gamma$ as $(\mu_i,u_i)_{i=1}^m$,
and define
\begin{equation}
 \rho_i=\mu_i/\mu_1,\qquad p_i=|u_i^Ty|^2/\|y\|^2,\qquad
 v_0=\mathbf1/\sqrt m.
\end{equation}
The weights sum to one, including energy in the exact nullspace. The symbol
$\lambda=\gamma h$ continues to denote relative bandwidth, not an eigenvalue.

Let $Q$ have orthonormal columns spanning $v_0^\perp$, and let
$t_\gamma(c)=(\tanh(\gamma(z_a-c)))_a$. Define
\begin{equation}
 H_\gamma=\frac1{hm}\int_{\mathbb R}
 Q^Tt_\gamma(c)t_\gamma(c)^TQ\,dc
 =-\frac2{hm}Q^T[d_{ab}\coth(\gamma d_{ab})]_{a,b}Q,
 \quad d_{ab}=z_a-z_b,
 \label{eq:app-slope-integral}
\end{equation}
where the diagonal value inside brackets is $1/\gamma$. Retain a finite set
$\mathcal F$ of absent exterior lattice centers such that the remaining absent
centers form two tails, starting at $c_L<z_{\min}$ and $c_R>z_{\max}$. Put
\begin{equation}
 T_\gamma=\frac1m\sum_{c\in\mathcal F}
 Q^Tt_\gamma(c)t_\gamma(c)^TQ,\qquad S_\gamma=H_\gamma-T_\gamma,
 \label{eq:app-slope-corrected}
\end{equation}
and let $\beta_1\ge\cdots\ge\beta_{m-1}$ be the eigenvalues of $S_\gamma$.
For any $0<\vartheta<\pi/2$, define
\begin{equation}
 \delta_\gamma=
 \frac{8\gamma\|z-\bar z\mathbf1\|^2\sec^4\vartheta}
 {3hm[\exp(2\pi\vartheta/(\gamma h))-1]},\qquad
 \bar z=m^{-1}\sum_a z_a,\qquad \ell_\gamma=v_0^TK_\gamma v_0.
 \label{eq:app-slope-lattice-allowance}
\end{equation}

\paragraph{Output-error lower bound.}
Set $\bar\rho_1=1$ and
\begin{equation}
 \bar\rho_i(\gamma)=\min\left\{1,
 \frac{\beta_{i-1}(\gamma)+\delta_\gamma}{\ell_\gamma}\right\},
 \qquad 2\le i\le m.
 \label{eq:app-slope-upper-rates}
\end{equation}
Then $\rho_i\le\bar\rho_i\le1$. Gradient descent from $w_0=0$ on
$\frac12\|J_\gamma w-y\|^2$, with $\eta=1/(2\mu_1)$, satisfies
\begin{equation}
 \frac{\|J_\gamma w_n-y\|^2}{\|y\|^2}
 \ge\sum_{i=1}^m p_i(\gamma)
 \left(1-\frac{\bar\rho_i(\gamma)}2\right)^{2n}
 \qquad(n\ge0).
 \label{eq:app-slope-output-bound}
\end{equation}
Known exact zero eigenvalues may have their endpoints set to zero; in
particular, $\operatorname{rank}K_\gamma\le W+1$. The conclusion depends on
target energy in the bounded slow directions; it does not assert that every
target is slow at small slope.

\subsection{The continuous-center kernel exposes gamma}
Removing the constant sample pattern makes the integral finite because
$Q^Tt_\gamma(c)$ decays exponentially at both ends. The identity
$1-\tanh u\tanh v=\coth(u-v)(\tanh u-\tanh v)$ gives
\begin{equation}
 \int_{\mathbb R}[1-\tanh(\gamma(z_a-c))\tanh(\gamma(z_b-c))]\,dc
 =2d_{ab}\coth(\gamma d_{ab}),
\end{equation}
with diagonal limit $2/\gamma$. Multiplication by $Q^T$ and $Q$ cancels the
constant term and proves Eq.~\eqref{eq:app-slope-integral}.

Write $s_\gamma(t)=\frac\gamma2\operatorname{sech}^2(\gamma t)$. The identity
$\tanh(\gamma\,\cdot)=s_\gamma*\operatorname{sign}$ follows because both sides
vanish at zero and have derivative $2s_\gamma$. For the convention
$\widehat g(\omega)=\int g(t)e^{-i\omega t}\,dt$, substitution
$r=e^{2\gamma t}$, the beta integral, and the gamma-function reflection formula
yield
\begin{align}
 \widehat s_\gamma(\omega)
 &=\int_0^\infty\frac{r^{-i\omega/(2\gamma)}}{(1+r)^2}\,dr
 =\Gamma\!\left(1-\frac{i\omega}{2\gamma}\right)
  \Gamma\!\left(1+\frac{i\omega}{2\gamma}\right)\notag\\
 &=\frac{\pi\omega/(2\gamma)}{\sinh(\pi\omega/(2\gamma))}
 =:M_\gamma(\omega),\qquad M_\gamma(0)=1.
\end{align}
For $u=Qv$, the step combination
$g_0(c)=\sum_a u_a\operatorname{sign}(c-z_a)$ has compact support since
$\sum_a u_a=0$. Its distributional derivative is
$2\sum_a u_a\delta_{z_a}$, so
$\widehat g_0(\omega)=2\sum_a u_a e^{-i\omega z_a}/(i\omega)$.
Parseval therefore gives
\begin{equation}
 v^TH_\gamma v=\frac2{\pi hm}\int_{\mathbb R}
 \frac{M_\gamma(\omega)^2}{\omega^2}
 \left|\sum_a(Qv)_a e^{-i\omega z_a}\right|^2\,d\omega.
 \label{eq:app-slope-fourier}
\end{equation}
The zero-frequency singularity is removable because the numerator vanishes
there. The direction's coefficients $(Qv)_a$ determine its frequency content;
gamma changes its weight through $M_\gamma^2$. Since
$d\log(z/\sinh z)/dz=1/z-\coth z<0$ for $z>0$, increasing gamma increases
$M_\gamma$ at every fixed frequency. Thus $H_\Gamma\succeq H_\gamma$ for
$\Gamma\ge\gamma$, and min-max gives ordered growth of the integral
eigenvalues. The eigenvectors may change. Exterior-center subtraction and
normalization by the largest finite eigenvalue need not preserve this order.

\subsection{Transfer to a finite center lattice}
For $u\perp\mathbf1$, let
$g(\zeta)=\sum_a u_a\tanh(\gamma(\zeta-z_a))$, analytic in
$|\operatorname{Im}\zeta|<\pi/(2\gamma)$. Subtract the common translate at
$\bar z$, write each difference as an integral of its derivative, and apply
Minkowski and Cauchy--Schwarz. On the strip boundaries this gives
\begin{equation}
 \int_{\mathbb R}|g(t\pm i\vartheta/\gamma)|^2\,dt
 \le\frac{4\gamma}{3}\sec^4\vartheta
 \|z-\bar z\mathbf1\|^2\|u\|^2.
 \label{eq:app-slope-strip}
\end{equation}
Indeed,
$|\operatorname{sech}^2(t+i\vartheta)|\le
\sec^2\vartheta\operatorname{sech}^2t$ and
$\int\operatorname{sech}^4t\,dt=4/3$, so the $L^2$ norm of a translate
difference is at most
$|z_a-\bar z|\sqrt{4\gamma/3}\sec^2\vartheta$.

Shift the contour for the analytic function $g(\zeta)^2$, not $|g(\zeta)|^2$.
Its transform is bounded by Eq.~\eqref{eq:app-slope-strip} times
$e^{-\vartheta|\omega|/\gamma}$. Poisson summation on $c_0+h\mathbb Z$
and the geometric sum over frequencies $2\pi k/h$, $k\ne0$, imply
\begin{equation}
 \left\|\frac1m\sum_{c\in c_0+h\mathbb Z}
 Q^Tt_\gamma(c)t_\gamma(c)^TQ-H_\gamma\right\|
 \le\delta_\gamma.
\end{equation}
The lattice origin enters only through phases of modulus one.

The finite kernel removes all absent centers. Let $T_{\rm tail}$ contain the
absent centers not retained in $\mathcal F$. Since $Q^T\mathbf1=0$ and
$1-\tanh t\le2e^{-2t}$ for $t\ge0$, a right-tail projected outer product
divided by $m$ has norm at most $4e^{-4\gamma(c-z_{\max})}$; the left tail
is analogous. Hence
\begin{equation}
 0\preceq T_{\rm tail}\preceq\tau_\gamma I,\qquad
 \tau_\gamma=\frac{4[e^{-4\gamma(c_R-z_{\max})}
 +e^{-4\gamma(z_{\min}-c_L)}]}{1-e^{-4\gamma h}}.
\end{equation}
Writing the lattice discrepancy as $R_{\rm lattice}$ yields
\begin{equation}
 Q^TK_\gamma Q=S_\gamma+R_{\rm lattice}-T_{\rm tail},\qquad
 \|R_{\rm lattice}\|\le\delta_\gamma,
\end{equation}
and consequently
\begin{equation}
 S_\gamma-(\delta_\gamma+\tau_\gamma)I
 \preceq Q^TK_\gamma Q\preceq S_\gamma+\delta_\gamma I.
 \label{eq:app-slope-enclosure}
\end{equation}
The bias vanishes under $Q$. The unretained exterior tail can only lower the
finite kernel; it is unnecessary for its upper eigenvalue bound.

\subsection{Relative rates and GD output error}
If $\kappa_i$ are the descending eigenvalues of $Q^TK_\gamma Q$, min-max and
codimension-one interlacing give
\begin{equation}
 \beta_i-\delta_\gamma-\tau_\gamma\le\kappa_i\le\beta_i+\delta_\gamma,
 \qquad \mu_i\ge\kappa_i\ge\mu_{i+1}.
\end{equation}
Thus $\mu_i\le\beta_{i-1}+\delta_\gamma$ for $i\ge2$. In exact arithmetic
this upper endpoint is nonnegative. The Rayleigh quotient
$\ell_\gamma=\|J_\gamma^Tv_0\|^2\le\mu_1$ is at least one because the bias
is present. Dividing by it and imposing $\rho_i\le1$ proves
Eq.~\eqref{eq:app-slope-upper-rates}.

The residual $r_n=J_\gamma w_n-y$ satisfies
\begin{equation}
 r_{n+1}=(I-\eta K_\gamma)r_n,\qquad r_0=-y,\qquad
 \frac{\|r_n\|^2}{\|y\|^2}=
 \sum_i p_i(1-\eta\mu_i)^{2n}.
\end{equation}
At $\eta=1/(2\mu_1)$ the factors are $(1-\rho_i/2)^{2n}$, decreasing in
$\rho_i\in[0,1]$. Replacing $\rho_i$ by its upper bound and summing against
$p_i\ge0$ proves Eq.~\eqref{eq:app-slope-output-bound}. This step is a fixed
conservative fraction of the stability limit $2/\mu_1$, not the maximal stable
step. For any nonoscillatory step $0<\eta\mu_1\le1$, the proof replaces
$\rho/2$ by $\eta\mu_1\rho$.

Only energy in the exact nullspace contributes a permanent floor
$p_0=\sum_{\mu_i=0}p_i$. A positive mode is representable because
$u_i=J_\gamma(J_\gamma^Tu_i/\mu_i)$, even when its decay is slow. The weights
$p_i$ use the actual finite-kernel eigenvectors; replacing them with
continuous-center eigenvectors would require another argument.

\subsection{Numerical evaluation and sources of slack}
The analytic statement uses exact $H_\gamma$ and $S_\gamma$. Numerical
evaluation uses a positive center-quadrature factor $A$ and an exterior factor
$C$, forming $AA^T-CC^T$. It does not truncate a Fourier series: quadrature
evaluates the same integral whose Fourier identity explains the mechanism.

Truncating the integral at
$\pm(\max_a|z_a|+P/\gamma)$ omits a positive matrix with norm at most
$2e^{-4P}/(h\gamma)$, added to the numerical upper endpoint. The exterior tail
enters the lower side of the kernel enclosure. For $Z=[A,C]$, its truncated SVD $Z_r$, and
$D=\operatorname{diag}(I,-I)$, expansion of $Z=Z_r+E$ gives
\begin{equation}
 \|ZDZ^T-Z_rDZ_r^T\|
 \le2\sigma_1(Z)\sigma_{r+1}(Z)+\sigma_{r+1}(Z)^2.
\end{equation}
This allowance enters both endpoints. Pad the compressed spectrum with zeros
between its positive and negative eigenvalues to recover dimension $m-1$.

Quadrature and arithmetic error are separate from these analytic tails.
Quadrature refinement and sensitivity checks support a numerical evaluation;
they are not rigorous error certificates. A multiple of machine epsilon is an
empirical guard unless justified by an arithmetic-error proof. Clipping a
negative computed upper endpoint at zero cannot repair a violated enclosure.
Computed curves should therefore be identified as checked floating-point
evaluations of the theorem.

Omitting numerically unresolved target energy weakens the lower curve
conservatively. Report unresolved spectral energy separately from the directly
computed projection residual; neither alone proves an exact nonrepresentability
floor in finite precision. The analytic bound also loses sharpness through the
interlacing index shift and the use of $\ell_\gamma$ instead of $\mu_1$.
Additive lattice and numerical allowances can dominate very small eigenvalues.
Retaining more exterior centers or improving arithmetic reduces some of this
slack, but does not remove the interlacing or normalization gaps. These are
distinct from inaccurate numerical evaluation.

All conclusions concern fixed features and the specified sample norm.
Population error, evolving centers or slopes, and Adam require additional
analysis or empirical evidence.
